\documentclass{article}
\usepackage{graphicx} % Required for inserting images
\usepackage{amssymb}
\usepackage{tikz}
\usepackage{tabularx}
\usepackage[czech]{babel}
\usepackage{csquotes}
\usepackage{hyperref}
\usepackage{gensymb}
\usepackage[thinc]{esdiff}
\hypersetup{
    colorlinks=true,
    linkcolor=blue,
    filecolor=magenta,      
    urlcolor=cyan,
    pdftitle={Calculus 2 Knowledge Base},
}

\urlstyle{same}
\usepackage{amsmath}
\usepackage{centernot}
\DeclareQuoteAlias{german}{czech}
\MakeOuterQuote{"}

\usepackage{booktabs,xltabular}
\usepackage[a4paper, total={6in, 8in}]{geometry}
\graphicspath{ {./images/} }
\newcommand*{\captionsource}[2]{%
  \caption[{#1}]{%
    #1%
    \\\hspace{\linewidth}%
    \textbf{Zdroj:} #2%
  }%
}
\title{Calculus 2 Knowledge Base}
\author{Jakub \textit{\{knedl1k\}} Adamec}

\newcommand*{\vv}[1]{\vec{\mkern0mu#1}}

\begin{document}

\pagenumbering{gobble}
\maketitle
\newpage
\pagenumbering{arabic}

\section{Řady}

Co se počítá: 
\begin{enumerate}
    \item určit, zda řada (absolutně) konverguje, nebo diverguje.
    \item poloměr konvergence.
    \item součet řady.
\end{enumerate}
\\
(Limitní) \textbf{Kritéria konvergence} pro řadu $\sum_{n=0}^{\infty} c_n$:
\\
(i) podílové:

    \begin{itemize}
        \item $\lim_{n \to \infty} \frac{|c_{n +1}|}{|c_n|} < 1$ ... řada (absolutně) konverguje.
        \item $\lim_{n \to \infty} \frac{|c_{n +1}|}{|c_n|} > 1$ ... řada diverguje.
        \item $\lim_{n \to \infty} \frac{|c_{n +1}|}{|c_n|} = 1$ ... nic nám neříká, zkus jiné kritérium.
    \end{itemize}
(ii) odmocninové:

    \begin{itemize}
        \item $\lim_{n \to \infty} \sqrt[n]{|c_n|} < 1$ ... řada (absolutně) konverguje.
        \item $\lim_{n \to \infty} \sqrt[n]{|c_n|} > 1$ ... řada diverguje.
        \item $\lim_{n \to \infty} \sqrt[n]{|c_n|} = 1$ ... nic nám neříká, zkus jiné kritérium.
    \end{itemize}
(iii) srovnávací:

Najdeme vhodnou řadu $\sum_{n=0}^{\infty} d_n$, s jejíž pomocí určíme konvergenci původní $\sum_{n=0}^{\infty} c_n$.
    \begin{itemize}
        \item $0 < \lim_{n \to \infty} \frac{c_n}{d_n} < \infty$ ... $\sum_{n=0}^{\infty} c_n \text{konverguje} \iff \sum_{n=0}^{\infty} d_n \text{konverguje}$.
    \end{itemize}
\textbf{Poloměr konvergence} $R$ řady $\sum_{n=0}^{\infty} a_n (x-x_0)^n$ je hodnota $0 \leq R \leq +\infty$. 

Řada $\sum_{n=0}^{\infty} a_n (x-x_0)^n$ se nazývá mocninná řada se středem v $x_0$ pro nějaké $x_0 \in \mathbb{R}$ a všechna $a_n \in \mathbb{R}$. Pro každé $x \in \mathbb{R}$ platí:
\begin{itemize}
    \item $|x-x_0| < R$ ... řada konverguje absolutně.
    \item $|x-x_0| > R$ ... řada diverguje. 
\end{itemize}
Jestliže $R = + \infty$, pak řada konverguje na celé přímce $\mathbb{R}$ a jestliže $0 < R < + \infty$, pak řada určitě konverguje na intervalu $(x_0 - R, x_0 + R)$ a nemusí konvergovat i v krajních bodech tohoto intervalu.
\\ \\
Příklady na další straně.
\newpage

\textbf{Příklad} Vyšetřete konvergenci řady $\sum_{n=0}^{\infty} \frac{n^2}{3^n}$.

\[
\lim_{n \to \infty} \frac{|c_{n+1}|}{|c_n|} \implies \frac{(n+1)^2}{3^{n+1}} \cdot \frac{3^n}{n^2} = \frac{(n+1)^2}{3 n^2} \xrightarrow{n \to \infty} \frac{1}{3} < 1 \text{ ... řada konverguje.}
\] 

\textbf{Příklad} Určete poloměr konvergence a sečtěte mocninnou řadu $\sum_{n=0}^{\infty} (n + (-2)^n) x^n$ na vnitřku oboru konvergence.

1. podílové kritérium




\newpage
\section{Fourierova řada}
\textbf{Příklad} Nalezněte Fourierovu řadu pro periodické rozšíření funkce

\[
f(t) = \begin{cases} \cos t  & \text{if } t \in [- \frac{\pi}{2}, \frac{\pi}{2}), \\
                      0      & \text{if } t \in [\frac{\pi}{2}, \frac{3 \pi}{2}).      %
        \end{cases}
\]

\begin{enumerate}
    \item Určit periodu $T$ a úhlovou frekvenci $\omega$.\\
    $T=2 \pi$, $\omega = \frac{2 \pi}{T} = \frac{2 \pi}{2 \pi} = 1$.
    \item \textit{Optional.} Zhodnotit sudost/lichost funkce, tím vyřadit část řešení.
    \item Spočítat $a_0$.\\
    \[
    a_0 = \frac{2}{T} \int_{-\frac{\pi}{2}}^{\frac{3\pi}{2}} f(t) \cdot cos(k \omega t) dt = \int_{-\frac{\pi}{2}}^{\frac{\pi}{2}} cos(t) \cdot cos(0t) dt = \int_{-\frac{\pi}{2}}^{\frac{\pi}{2}} cos(t) dt = [sin t]_{-\frac{\pi}{2}}^{\frac{\pi}{2}} = 0
    \]
    \item Spočítat $a_k$.
    \[
    a_k = \frac{2}{T} \int_{-\frac{\pi}{2}}^{\frac{3\pi}{2}} f(t) \cdot cos(k \omega t) dt = \int_{-\frac{\pi}{2}}^{\frac{\pi}{2}} cos(t) \cdot cos(kt) dt = \int_{-\frac{\pi}{2}}^{\frac{\pi}{2}} 
    \]
    
\end{enumerate}

\section{Limity}


\newpage
\section{Parciální derivace}

Co se zde může počítat: 
\begin{enumerate}
    \item derivace vůči nějaké proměnné.
    \item spojitost funkce, spojitost parciální derivace (přes limity).
    \item totální diferenciál.
    \item tečná rovina.
    \item derivace ve směru.
    \item úhel, který svírá tečná rovina s jinou rovinou.
\end{enumerate}
\textbf{Příklad} Pro funkci $f(x, y) = \ln{x y^2}$ najděte parciální derivace $\diffp{f}{x}$ a $\diffp{f}{y}$ a obory jejich existence.

    V bodě $a_0 = (1, 1)$ určete totální diferenciál, tečnou rovinu a derivaci ve směru $\vec{u} = (\frac{\sqrt{2}}{2}, -\frac{\sqrt{2}}{2})$.

    Určete úhel, který svírá tečná rovina se základnou (tj. s rovinou $z=0$).
    \begin{enumerate}
        \item $\diffp{f}{x} = \frac{1}{x} $, $\diffp{f}{y} = \frac{2y}{x y}$.
            $x > 0, y \neq 0$. \\ (Podmínky se určují vždy z průniku podmínek funkce a parciální derivace.)
        \item $\nabla f(1, 1) = (1, 2)$.
        \item $z = f(a_0) + \diffp{f}{x}(a_0) \cdot (x-x_0) + \diffp{f}{y}(a_0) \cdot (y-y_0)= 0 + 1(x-1) + 2(y-1) = x+2y-3$. \\
        Jiný způsob výpočtu: $z = f(a_0) + 
        \begin{pmatrix}
            \diffp{f}{x}(a_0), \diffp{f}{y}(a_0)
        \end{pmatrix}
        \cdot 
        \begin{pmatrix}
            x-x_0 \\           
            y-y_0
        \end{pmatrix}
        = x+2y-3$.\\
        (Za $x_0$ a $y_0$ dosadím bod $a_0$, zapamatuješ si to díky indexům, nulu dosadit do nuly.)
        \item $\diffp{f}{u}(a_0) = (1, 2) \cdot 
        \begin{pmatrix}
            \frac{\sqrt{2}}{2} \\           
            -\frac{\sqrt{2}}{2}
        \end{pmatrix} 
        = \frac{\sqrt{2}}{2} - \sqrt{2} = -\frac{\sqrt{2}}{2}$. \\
        (Vezmu totální diferenciál v bodě a vynásobím ho směrem $\vec{u}$. To je vše.)
        \item Díky zadání znám první normálový vektor $\vec{n_1} = (0, 0, 1)$. Druhý pak určim z rovnice tečné roviniy z bodu 3. $\vec{n_2} = (\nabla f(a_0), -1) = (1, 2, -1)$. \\
        $\cos{\alpha} = \frac{|\vec{n_1} \cdot \vec{n_2}|}{||{\vec{n_1}|| \cdot ||\vec{n_2}||}}$, $\alpha \in [0, \frac{\pi}{2} ]$. $\implies \alpha = \arccos{\frac{|\vec{n_1} \cdot \vec{n_2}|}{||{\vec{n_1}|| \cdot ||\vec{n_2}||}}} = \arccos{\frac{1}{1 \cdot \sqrt{6}}} \approx 65,9 \degree$.
    \end{enumerate}
    Více příkladů \href{https://math.fel.cvut.cz/en/people/korbemir/mat2/2022z/4.cviceni_mat2_2022_z.pdf}{zde}, případně \href{https://math.fel.cvut.cz/en/people/korbemir/mat2/2024z/5.cviceni_mat2_2024_z.pdf}{zde}.
    













\end{document}